\documentclass[a4paper,11pt]{article}

\usepackage{revy}
\usepackage[utf8]{inputenc}
\usepackage[T1]{fontenc}
\usepackage[danish]{babel}


\revyname{MatematikRevyen}
\revyyear{2026}
\version{1}
\eta{$3$ minutter}
\status{Færdig}

\title{Brand}
\author{Julius Pallesen '24}
\melody{Monroe:``Regarde !''}

\begin{document}
\maketitle

\begin{roles}
\role{S}[Sanger] Vokalstærk operasanger
\role{B1}[Brandmand 1] Brandmænd (M/K) der kommer ind lige inden omkv. 2
\role{B2}[Brandmand 2]
\end{roles}

\begin{song}
    \sings{S}[Intro] Sidder til tim' – keder mig lidt
Så jeg ta'r en lur – for jeg er lidt træt

\sings{S}[Bro] Det ik' for – at være sur
Men jeg – syn's ik' de dur
At en – brandalarm ødelægger min lur
\scene{Musikken afbydes kort og en brandarlarm lyder (sørg for at den ikke kan forveksles med en rigtig brandalarm)}

\sings{S}[Omkv] Der brand igen, der brand igen!
Skal jeg dog bare, tage turen hjem?
Der brand igen, der brand igen!
Så er det nu – at vi skal brug' – nødudgangen

\act{Operastykke}

\sings{S}[Vers] Er det mon i A103
at et lille bål er blevet tændt af
Eller er det bare fysik
Der laver forsøg i pyroteknik
Eller måske er det kemi
Der laver flammer man ikke kan se
Kunne det være mig, der har glemt
At slukke ovnen i S01?

\sings{S}[Bro] Det ik' for – at være sur
Men jeg – syn's ik' de dur
At ovnen – ikke har et ur

\scene{Musikken afbrydes igen kort. Denne gang af udrykning. To brandmænd træder ind på scenen og beklager sig.}

\sings{B1} Hvad har i gang i denne gang? Vi kommer fanme rendende hver uge.

\scene{Musikken går i gang igen.}

\sings{S}[Omkv] \act{Til brandmændene}
Der brand igen, der brand igen!
Nu skal i ind, og redde min ven!
Der brand igen, der brand igen!
Vi kommer aldrig til at se – HC(Ø) igen

\sings{B1}Der er ik' ægte brand
\sings{S}Shhh
\sings{B2}Det er falsk alarm
\sings{S}Shhh
Det er løgn, løgn, løgn

\act{Operastykke}

\sings{S}[Outro] Der ikke brand, der ikke brand
Så nu kan vi end'lig gå igang
\scene{Musikken slutter brat. Lys ned}

\end{song}

\end{document}

